\documentclass[11pt, a4paper]{ctexart}

\usepackage{assignment_style}

\begin{document}

\begin{center}
    \textbf{最小二乘 $m$ 次拟合多项式}
\end{center}

设
\[
    P_m(x)=\sum_{k=0}^{m}a_kx^k,
\]
则系数向量 $\bm a=(a_0,\ldots,a_m)^{\mathrm T}$ 满足正规方程组
\[
    \begin{pmatrix}
        \sum x_i^0 & \sum x_i^1 & \cdots & \sum x_i^m \\
        \sum x_i^1 & \sum x_i^2 & \cdots & \sum x_i^{m+1} \\
        \vdots & \vdots & \ddots & \vdots \\
        \sum x_i^m & \sum x_i^{m+1} & \cdots & \sum x_i^{2m}
    \end{pmatrix}
    \bm a
    =
    \begin{pmatrix}
        \sum y_ix_i^0 \\
        \sum y_ix_i^1 \\
        \vdots \\
        \sum y_ix_i^m
    \end{pmatrix}.
\]

% =========================================
% 1
% =========================================
\begin{problem}{1. 已知数据表如下：}
    \[
        \begin{array}{c|cccccccc}
            x_i & 1.36 & 1.49 & 1.73 & 1.81 & 1.95 & 2.16 & 2.28 & 2.48 \\ \hline
            y_i & 14.094 & 15.069 & 16.844 & 17.378 & 18.435 & 19.949 & 20.963 & 22.495
        \end{array}
    \]
    求其一次最小二乘拟合多项式。
\end{problem}

\begin{solution}
    取拟合多项式
    \[
        P_1(x)=a_0+a_1x.
    \]
    计算各和式：
    \[
        \sum x_i^0=8,\qquad \sum x_i=15.26,\qquad \sum x_i^2=30.1556,
    \]
    \[
        \sum y_i=145.227,\qquad \sum x_iy_i=284.8363.
    \]
    因而正规方程为
    \[
        \begin{pmatrix}
            8 & 15.26 \\
            15.26 & 30.1556
        \end{pmatrix}
        \begin{pmatrix}
            a_0 \\
            a_1
        \end{pmatrix}
        =
        \begin{pmatrix}
            145.227 \\
            284.8363
        \end{pmatrix}.
    \]
    解得
    \[
        a_0=3.9161,\qquad a_1=7.4639.
    \]
    所以一次最小二乘拟合多项式为
    \[
        P_1(x)=3.9161+7.4639x.
    \]
\end{solution}


% =========================================
% 2
% =========================================
\begin{problem}{2. 已知数据表如下：}
    \[
        \begin{array}{c|ccccccccc}
            x_i & 1 & 3 & 4 & 5 & 6 & 7 & 8 & 9 & 10 \\ \hline
            y_i & 10 & 5 & 4 & 2 & 1 & 1 & 2 & 3 & 4
        \end{array}
    \]
    求其二次最小二乘拟合多项式。
\end{problem}

\begin{solution}
    设
    \[
        P_2(x)=a_0+a_1x+a_2x^2.
    \]
    计算得到
    \[
        \sum x_i^0=9,\qquad \sum x_i=53,\qquad \sum x_i^2=381,
    \]
    \[
        \sum x_i^3=3017,\qquad \sum x_i^4=25317,
    \]
    \[
        \sum y_i=32,\qquad \sum x_iy_i=147,\qquad \sum x_i^2y_i=1025.
    \]
    故正规方程组为
    \[
        \begin{pmatrix}
            9 & 53 & 381 \\
            53 & 381 & 3017 \\
            381 & 3017 & 25317
        \end{pmatrix}
        \begin{pmatrix}
            a_0 \\
            a_1 \\
            a_2
        \end{pmatrix}
        =
        \begin{pmatrix}
            32 \\
            147 \\
            1025
        \end{pmatrix}.
    \]
    解得
    \[
        a_0=13.4597,\qquad a_1=-3.6053,\qquad a_2=0.2676.
    \]
    因而二次最小二乘拟合多项式为
    \[
        P_2(x)=13.4597-3.6053x+0.2676x^2.
    \]
\end{solution}


% =========================================
% 4
% =========================================
\begin{problem}{4. 用最小二乘法求一形如 $y=a+bx^2$（$a,b$ 为常数）的经验公式拟合下列数据：}
    \[
        \begin{array}{c|ccccc}
            x_i & 19 & 25 & 31 & 38 & 44 \\ \hline
            y_i & 19.0 & 32.3 & 49.0 & 73.3 & 97.8
        \end{array}
    \]
\end{problem}

\begin{solution}
    令
    \[
        X_i=x_i^2,
    \]
    则问题化为对线性模型
    \[
        y=a+bX
    \]
    作一次最小二乘拟合。此时
    \[
        \sum X_i^0=5,\qquad \sum X_i=5327,\qquad \sum X_i^2=7277699,
    \]
    \[
        \sum y_i=271.4,\qquad \sum X_iy_i=369320.
    \]
    因而正规方程为
    \[
        \begin{pmatrix}
            5 & 5327 \\
            5327 & 7277699
        \end{pmatrix}
        \begin{pmatrix}
            a \\
            b
        \end{pmatrix}
        =
        \begin{pmatrix}
            271.4 \\
            369320
        \end{pmatrix}.
    \]
    解得
    \[
        a=0.9726,\qquad b=0.0500.
    \]
    所求经验公式为
    \[
        y=0.9726+0.0500x^2.
    \]
\end{solution}


% =========================================
% 6
% =========================================
\begin{problem}{6. 用最小二乘法求一形如 $y=ae^{bx}$ 的经验公式拟合下列数据：}
    \[
        \begin{array}{c|cccccccc}
            x_i & 1 & 2 & 3 & 4 & 5 & 6 & 7 & 8 \\ \hline
            y_i & 15.3 & 20.5 & 27.4 & 36.6 & 49.1 & 65.6 & 87.87 & 117.6
        \end{array}
    \]
\end{problem}

\begin{solution}
    对模型
    \[
        y=ae^{bx}
    \]
    两边取对数，得
    \[
        \ln y=\ln a+bx.
    \]
    令
    \[
        Y_i=\ln y_i,\qquad c=\ln a,
    \]
    则转化为线性模型
    \[
        Y=c+bx.
    \]
    由数据得
    \[
        \sum x_i^0=8,\qquad \sum x_i=36,\qquad \sum x_i^2=204,
    \]
    \[
        \sum Y_i=29.9795,\qquad \sum x_iY_i=147.1406.
    \]
    故正规方程为
    \[
        \begin{pmatrix}
            8 & 36 \\
            36 & 204
        \end{pmatrix}
        \begin{pmatrix}
            c \\
            b
        \end{pmatrix}
        =
        \begin{pmatrix}
            29.9795 \\
            147.1406
        \end{pmatrix}.
    \]
    解得
    \[
        c=2.4367,\qquad b=0.2913.
    \]
    因而
    \[
        a=e^c=e^{2.4367}\approx 11.4358.
    \]
    所求经验公式为
    \[
        y=11.4358e^{0.2913x}.
    \]
\end{solution}


% =========================================
% 7
% =========================================
\begin{problem}{7. 某化学反应中，测得生成物的质量浓度 $y$ 与时间 $t$ 的关系见下表：}
    \[
        \begin{array}{c|cccccccccc}
            t_i & 1 & 2 & 3 & 4 & 6 & 8 & 10 & 12 & 14 & 16 \\ \hline
            y_i & 4.00 & 6.41 & 8.01 & 8.79 & 9.53 & 9.86 & 10.33 & 10.42 & 10.53 & 10.61
        \end{array}
    \]
    用最小二乘法求一形如
    \[
        y=\frac{t}{at+b}
    \]
    的经验公式。
\end{problem}

\begin{solution}
    将模型改写为
    \[
        \frac{1}{y}=a+\frac{b}{t}.
    \]
    令
    \[
        X_i=\frac{1}{t_i},\qquad Y_i=\frac{1}{y_i},
    \]
    则问题转化为线性拟合
    \[
        Y=a+bX.
    \]
    由实验数据可得
    \[
        \sum X_i^0=10,\qquad \sum X_i=2.6923,\qquad \sum X_i^2=1.4930,
    \]
    \[
        \sum Y_i=1.2330,\qquad \sum X_iY_i=0.4586.
    \]
    因而正规方程为
    \[
        \begin{pmatrix}
            10 & 2.6923 \\
            2.6923 & 1.4930
        \end{pmatrix}
        \begin{pmatrix}
            a \\
            b
        \end{pmatrix}
        =
        \begin{pmatrix}
            1.2330 \\
            0.4586
        \end{pmatrix}.
    \]
    解得
    \[
        a=0.0789,\qquad b=0.1649.
    \]
    所求经验公式为
    \[
        y=\frac{t}{0.0789t+0.1649}.
    \]
\end{solution}


% =========================================
% 10
% =========================================
\begin{problem}{10. 证明：函数系 $\{\varphi_0,\ldots,\varphi_n\}$ 线性无关的充要条件是 Gram 行列式}
    \[
        G(\varphi_0,\ldots,\varphi_n)\neq 0.
    \]
\end{problem}

\begin{myproof}
    记 Gram 矩阵为
    \[
        G=\bigl((\varphi_i,\varphi_j)\bigr)_{i,j=0}^{n}.
    \]

    \textbf{必要性} 若 $\{\varphi_0,\ldots,\varphi_n\}$ 线性无关，则对任意非零向量
    \[
        \bm a=(a_0,\ldots,a_n)^{\mathrm T}\neq \bm 0,
    \]
    定义
    \[
        \varphi=\sum_{i=0}^{n}a_i\varphi_i.
    \]
    由线性无关性知 $\varphi\neq 0$，于是
    \[
        (\varphi,\varphi)>0.
    \]
    另一方面，
    \[
        (\varphi,\varphi)=\sum_{i=0}^{n}\sum_{j=0}^{n}a_i a_j(\varphi_i,\varphi_j)
        =\bm a^{\mathrm T}G\bm a.
    \]
    因而对任意非零 $\bm a$ 有
    \[
        \bm a^{\mathrm T}G\bm a>0,
    \]
    即 $G$ 正定，从而
    \[
        \det G\neq 0.
    \]

    \textbf{充分性} 若 $\det G\neq 0$，设
    \[
        \sum_{i=0}^{n}a_i\varphi_i=0.
    \]
    将上式分别与 $\varphi_j\ (j=0,1,\ldots,n)$ 作内积，得
    \[
        \sum_{i=0}^{n}a_i(\varphi_i,\varphi_j)=0,\qquad j=0,1,\ldots,n.
    \]
    这就是齐次线性方程组
    \[
        G\bm a=0.
    \]
    由于 $\det G\neq 0$，该方程组只有零解，即
    \[
        \bm a=0.
    \]
    故 $\{\varphi_0,\ldots,\varphi_n\}$ 线性无关。

    综上，函数系 $\{\varphi_0,\ldots,\varphi_n\}$ 线性无关当且仅当 $G(\varphi_0,\ldots,\varphi_n)\neq 0$。
\end{myproof}


% =========================================
% 12
% =========================================
\begin{problem}{12. 求 $f(x)=\cos x,\ x\in\left[-\dfrac{\pi}{2},\dfrac{\pi}{2}\right]$ 在}
    \[
        H=\operatorname{Span}\{1,x^2,x^4\}
    \]
    中的最佳平方逼近多项式。
\end{problem}

\begin{solution}
    取基函数
    \[
        \varphi_0(x)=1,\qquad \varphi_1(x)=x^2,\qquad \varphi_2(x)=x^4,
    \]
    设最佳平方逼近多项式为
    \[
        p(x)=a_0+a_1x^2+a_2x^4.
    \]
    在内积
    \[
        (f,g)=\int_{-\pi/2}^{\pi/2}f(x)g(x)\,\dif x
    \]
    下，系数满足正规方程
    \[
        A\bm a=\bm b,
    \]
    其中
    \[
        A=\bigl((\varphi_i,\varphi_j)\bigr)_{i,j=0}^{2},\qquad
        \bm b=\bigl((f,\varphi_0),(f,\varphi_1),(f,\varphi_2)\bigr)^{\mathrm T}.
    \]
    计算得
    \[
        A=
        \begin{pmatrix}
            3.1416 & 2.5839 & 3.8252 \\
            2.5839 & 3.8252 & 6.7417 \\
            3.8252 & 6.7417 & 12.9380
        \end{pmatrix},
        \qquad
        \bm b=
        \begin{pmatrix}
            2.0000 \\
            0.9348 \\
            0.9585
        \end{pmatrix}.
    \]
    解得
    \[
        a_0=0.9996,\qquad a_1=-0.4964,\qquad a_2=0.0372.
    \]
    所以最佳平方逼近多项式为
    \[
        p(x)=0.9996-0.4964x^2+0.0372x^4.
    \]
\end{solution}


% =========================================
% 13(1)
% =========================================
\begin{problem}{13.(1) 求函数 $f(x)=\sqrt{x}$ 在区间 $[0,1]$ 上关于 $f(x)$ 的一次最佳平方逼近多项式。}
\end{problem}

\begin{solution}
    设所求一次多项式为
    \[
        p_1(x)=a_0+a_1x.
    \]
    取基函数
    \[
        \varphi_0(x)=1,\qquad \varphi_1(x)=x.
    \]
    在内积
    \[
        (f,g)=\int_{0}^{1}f(x)g(x)\,\dif x
    \]
    下，正规方程为
    \[
        \begin{pmatrix}
            (\varphi_0,\varphi_0) & (\varphi_1,\varphi_0) \\
            (\varphi_0,\varphi_1) & (\varphi_1,\varphi_1)
        \end{pmatrix}
        \begin{pmatrix}
            a_0 \\
            a_1
        \end{pmatrix}
        =
        \begin{pmatrix}
            (f,\varphi_0) \\
            (f,\varphi_1)
        \end{pmatrix}.
    \]
    计算得
    \[
        A=
        \begin{pmatrix}
            1 & \dfrac{1}{2} \\
            \dfrac{1}{2} & \dfrac{1}{3}
        \end{pmatrix},
        \qquad
        \bm b=
        \begin{pmatrix}
            \displaystyle\int_{0}^{1}x^{1/2}\,\dif x \\
            \displaystyle\int_{0}^{1}x^{3/2}\,\dif x
        \end{pmatrix}
        =
        \begin{pmatrix}
            \dfrac{2}{3} \\
            \dfrac{2}{5}
        \end{pmatrix}.
    \]
    即
    \[
        \begin{pmatrix}
            1 & \dfrac{1}{2} \\
            \dfrac{1}{2} & \dfrac{1}{3}
        \end{pmatrix}
        \begin{pmatrix}
            a_0 \\
            a_1
        \end{pmatrix}
        =
        \begin{pmatrix}
            \dfrac{2}{3} \\
            \dfrac{2}{5}
        \end{pmatrix}.
    \]
    解得
    \[
        a_0=\frac{4}{15},\qquad a_1=\frac{4}{5}.
    \]
    因而一次最佳平方逼近多项式为
    \[
        p_1(x)=\frac{4}{15}+\frac{4}{5}x.
    \]
\end{solution}

\end{document}
